\documentclass[11pt,a4paper]{article}

% --- Codificación e idioma ---
\usepackage[T1]{fontenc}
\usepackage[spanish,es-nodecimaldot]{babel}

% --- Matemáticas ---
\usepackage{amsmath, amsthm, amssymb}

% --- Iconos (requerido por \faEdit en el entorno ejemplo) ---
\usepackage{fontawesome5}

% --- Colores y Tipografía ---
\usepackage{xcolor}
\usepackage{lmodern}
\renewcommand{\familydefault}{\sfdefault}

% --- Gráficos ---
\usepackage{tikz}
\usetikzlibrary{babel} 
\usetikzlibrary{patterns.meta}
\usepackage{pgfplots}
\pgfplotsset{compat=1.18, colormap/viridis}

% --- Cajas personalizadas ---
\usepackage[many]{tcolorbox}
\tcbuselibrary{skins, breakable, theorems}

% --- Títulos y Secciones ---
\usepackage{titlesec}
\usepackage{etoc}

% --- Formato decorativo para \section ---
\newcommand{\neonrule}[1]{%
    \begin{tikzpicture}[remember picture, overlay]
        \fill[#1] (0,0.3) rectangle (0.2,0.45);
        \fill[#1] (0.25,0.3) rectangle (0.45,0.45);
        \fill[#1] (0.5,0.3) rectangle (0.7,0.45);
        \draw[#1, line width=1pt] (1.2,0.375) -- ({\linewidth-0.4cm},0.375);
        \fill[#1] ({\linewidth-0.5cm}, 0.3) rectangle (\linewidth,0.45);
    \end{tikzpicture}%
}

\titleformat{\section}
  {\normalfont\Large\bfseries}{\thesection}{1em}{}
  [\neonrule{teal}]
\titlespacing*{\section}{0pt}{3.5ex plus 1ex minus .2ex}{2.3ex plus .2ex}

% --- Corrección de \portadaseccion ---
\newcommand{\portadaseccion}[1]{%
    \clearpage
    \thispagestyle{empty}
    \vspace*{2cm}
    \section{#1}
    \vspace{1.5cm}
    \begingroup
        \etocsettocstyle{\Large\bfseries Contenido del tema\par\vspace{0.5cm}\hrule\vspace{0.5cm}}{\vspace{0.5cm}\hrule}
        \etocsetnexttocdepth{subsubsection}
        \localtableofcontents
    \endgroup
    \clearpage
}

% --- Estilos de tcolorbox ---
\tcbset{
    futuristicbase/.style={
        enhanced,
        breakable,
        boxrule=0.5pt,
        arc=2pt,
        fonttitle=\bfseries
    }
}

% Entornos numerados
\newtcbtheorem[number within=section]{teorema}{Teorema}{
    futuristicbase,
    colback=blue!2!white,
    colframe=blue!85!black,
    fuzzy shadow={1.2mm}{-1.2mm}{0mm}{0.3mm}{blue!50!black,opacity=0.5},
    title style={left color=blue!90!black, right color=blue!60!black}
}{thm}

\newtcbtheorem[number within=section]{definicion}{Definición}{
    futuristicbase,
    colback=cyan!4!white,
    colframe=cyan!80!black,
    fuzzy shadow={1.2mm}{-1.2mm}{0mm}{0.3mm}{cyan!50!black,opacity=0.5},
    title style={left color=cyan!80!black, right color=cyan!50!black}
}{def}

\newtcbtheorem[number within=section]{proposicion}{Proposición}{
    futuristicbase,
    colback=blue!2!gray!5!white,
    colframe=blue!50!black,
    fuzzy shadow={1.2mm}{-1.2mm}{0mm}{0.3mm}{gray!50!black,opacity=0.5},
    title style={left color=blue!60!black, right color=blue!40!black}
}{prop}

\newtcbtheorem[number within=section]{observacion}{Observación}{
    futuristicbase,
    colback=yellow!5!white,
    colframe=yellow!60!black,
    boxrule=0.4pt,
    coltitle=black,
    fuzzy shadow={1.2mm}{-1.2mm}{0mm}{0.3mm}{yellow!50!black,opacity=0.5},
    title style={left color=yellow!40!white, right color=yellow!5!white}
}{obs}

\newtcbtheorem[number within=section]{nota}{Nota}{
    futuristicbase,
    colback=yellow!15!orange!5!white,
    colframe=yellow!50!orange!90!black,
    boxrule=0.4pt,
    coltitle=white,
    fuzzy shadow={1.2mm}{-1.2mm}{0mm}{0.3mm}{orange!50!black,opacity=0.5},
    title style={left color=yellow!40!orange!90!black, right color=yellow!60!orange!70!black}
}{not}

% Entorno corregido con \faEdit disponible
\newtcbtheorem[number within=section]{ejemplo}{\faEdit\ Ejemplo}{
    futuristicbase,
    colback=teal!5!white,
    colframe=teal!80!black,
    boxrule=0pt,
    borderline={1pt}{0pt}{teal, dashed},
    coltitle=black,
    colbacktitle=teal!15!white
}{ej}

% --- Referencias (deben cargarse al final) ---
\usepackage{hyperref}
\usepackage[capitalise, spanish]{cleveref}

\hypersetup{
    colorlinks        = true,
    linkcolor         = blue!75!black,
    citecolor         = yellow!60!black,
    urlcolor          = cyan!70!black,
    pdftitle          = {Apuntes},
    pdfborder         = {0 0 0},
    bookmarksnumbered = true,
}

\crefname{tcb@cnt@definicion}{definición}{definiciones}
\crefname{tcb@cnt@observacion}{observación}{observaciones}
\crefname{tcb@cnt@nota}{nota}{notas}
\crefname{tcb@cnt@teorema}{teorema}{teoremas}
\crefname{tcb@cnt@proposicion}{proposición}{proposiciones}
\crefname{tcb@cnt@ejemplo}{ejemplo}{ejemplos}

\Crefname{tcb@cnt@definicion}{Definición}{Definiciones}
\Crefname{tcb@cnt@observacion}{Observación}{Observaciones}
\Crefname{tcb@cnt@nota}{Nota}{Notas}
\Crefname{tcb@cnt@teorema}{Teorema}{Teoremas}
\Crefname{tcb@cnt@proposicion}{Proposición}{Proposiciones}
\Crefname{tcb@cnt@ejemplo}{Ejemplo}{Ejemplos}

% --- Cabeceras y pie de página ---
\usepackage{fancyhdr}
\setlength{\headheight}{15pt} % Ajustado para evitar warning de altura
\pagestyle{fancy}
\fancyhf{}
\renewcommand{\headrulewidth}{0.4pt}
\renewcommand{\footrulewidth}{0.4pt}
\fancyhead[L]{\small\textcolor{gray}{\leftmark}}
\fancyhead[R]{\small\textcolor{gray}{Apuntes}}
\fancyfoot[C]{\small\textcolor{gray}{\thepage}}

% --- Estilo global de página ---
\pagecolor{gray!5}
\color{black}

% --- Metadatos ---
\title{\Huge \textbf{Apuntes de Análisis de Fourier}}
\author{\Large Juan Manuel Flores de la Cruz}
\date{Curso 2026/27}